% !TEX root = ../main.tex
\begin{abstract}
LLM agents increasingly depend on external skills to perform medical tasks, yet existing benchmarks measure only end-to-end task success---not whether the agent correctly executes the skill specification it was given.
We introduce \medskillbench, the first large-scale skill-level execution benchmark for medical AI agents, comprising 1,462 skills spanning eight medical domains---from bioinformatics pipelines to clinical reasoning guides---evaluated in isolated sandboxes with a mandatory read-first protocol, eight-dimensional binary scoring, and a novel Specification Groundability Index (\textsc{SGI}) that decomposes execution quality into core fidelity and auxiliary quality components.
Our evaluation reveals a systematic \emph{read-but-not-use} (\rbu) gap: agents invoke skills at 96.2\% but complete specified tasks at only 74.3\%, a 22-percentage-point specification grounding failure that decomposes into a dominant semantic grounding component ($\delta_{\text{sem}} = 0.242$) and a negative execution completion component ($\delta_{\text{exec}} = -0.023$), indicating that agents often bypass skill specifications entirely via parametric knowledge.
Stratifying by skill interface type shows that guide-style skills (no executable tool) exhibit the largest gap (0.269), while tool-based skills show substantially smaller gaps (0.069--0.111), identifying description groundability---not environment dependency---as the primary bottleneck.
Skill-augmented evaluation on 12 medical QA benchmarks further demonstrates that providing matched skills improves accuracy by \todo{X} percentage points on average, with tool-based skills yielding substantially larger gains than guide-style descriptions.
These findings establish description groundability as the primary bottleneck for medical agent skills, quantify the practical value of skill specifications, and provide actionable diagnostic signal for future skill development.
\end{abstract}
